% from ega0-1-fr.tex
\oldpage[0\textsubscript{I}]{72}

mod.~$\mathfrak JB/\mathfrak J^2B$ forment un système de générateurs de la
$(A/\mathfrak J)$-algèbre $B/\mathfrak JB$~; on voit aussitôt que les $c_ia_j$
forment un système de générateurs de la $(A/\mathfrak J^2)$-algèbre
$B/\mathfrak J^2B$, donc a) implique b).

Reste à prouver que b) entraîne c). L'hypothèse entraîne que $B_1$ est un
anneau noethérien, et comme $B_1=B_2/\mathfrak J^2B_2$, on a
$\mathfrak J^2B_1=0$, donc
$\mathfrak JB_1=\mathfrak JB_1/\mathfrak J^2B_1$ est un $B_0$-module de type
fini. Les conditions de (7.2.7) sont donc vérifiées par le système projectif
$(B_n)$ et $B$ est un anneau $\mathfrak JB$-adique. Soit
$(c_i)_{1\leq i\leq r}$ un système fini d'éléments de $B$ dont les classes
mod.~$\mathfrak JB$ engendrent la $(A/\mathfrak J)$-algèbre
$B/\mathfrak JB$, et dont les combinaisons linéaires à coefficients dans
$\mathfrak J$ sont telles que leurs classes mod.~$\mathfrak J^2B$ engendrent
le $B_0$-module $\mathfrak JB/\mathfrak J^2B$. Il existe un
$A$-homomorphisme continu $u$ de $A'=A\{T_1,\ldots,T_r\}$ dans $B$ qui se
réduit à $\varphi$ dans $A$ et est tel que $u(T_i)=c_i$ pour $1\leq i\leq r$
(7.5.3)~; si nous prouvons que $u$ est \emph{surjectif}, c) sera établi, car de
$u(A')=B$ on déduira $u(\mathfrak J^nA')=\mathfrak J^nB$, autrement dit $u$
sera un morphisme strict d'anneaux topologiques et $B$ sera donc isomorphe à un
quotient de $A'$ par un idéal fermé. Or, comme $B$ est complet pour la topologie
$\mathfrak JB$-adique, il suffit ([I], p.~18-07) de montrer que
l'homomorphisme $\operatorname{grad}(A')\to\operatorname{grad}(B)$ déduit
canoniquement de $u$ pour les filtrations $\mathfrak J$-adiques sur $A'$ et
$B$, est surjectif. Mais par définition, les homomorphismes
$A'/\mathfrak JA'\to B/\mathfrak JB$ et
$\mathfrak JA'/\mathfrak J^2A'\to\mathfrak JB/\mathfrak J^2B$ déduits de $u$
sont surjectifs~; par récurrence sur $n$, on en déduit aussitôt qu'il en est de
même de $\mathfrak JA'/\mathfrak J^nA'\to\mathfrak JB/\mathfrak J^nB$, et
\emph{a fortiori} de
$\mathfrak J^nA'/\mathfrak J^{n+1}A'\to
\mathfrak J^nB/\mathfrak J^{n+1}B$, ce qui achève la démonstration.

\subsection{Anneaux complets de fractions}
\label{subsection:0.7.6-fr}

\begin{env}[7.6.1]
\label{0.7.6.1-fr}
Soient $A$ un anneau linéairement topologisé, $(\mathfrak J_\lambda)$ un
système fondamental de voisinages de 0 dans $A$ formé d'idéaux, $S$ une partie
multiplicative de $A$. Soit $u_\lambda$ l'homomorphisme canonique
$A\to A_\lambda=A/\mathfrak J_\lambda$, et pour
$\mathfrak J_\mu\subset\mathfrak J_\lambda$, soit $u_{\lambda\mu}$
l'homomorphisme canonique $A_\mu\to A_\lambda$. Posons
$S_\lambda=u_\lambda(S)$, de sorte que
$u_{\lambda\mu}(S_\mu)=S_\lambda$. Les $u_{\lambda\mu}$ donnent
canoniquement des homomorphismes surjectifs
$S_\mu^{-1}A_\mu\to S_\lambda^{-1}A_\lambda$, pour lesquels ces anneaux
forment un système projectif~; désignons par $A\{S^{-1}\}$ la limite
projective de ce système. Cette définition ne dépend pas du système fondamental
de voisinages $(\mathfrak J_\lambda)$ choisi~; en effet~:
\end{env}

\begin{proposition}[7.6.2]
\label{0.7.6.2-fr}
L'anneau $A\{S^{-1}\}$ est topologiquement isomorphe au séparé complété de
l'anneau $S^{-1}A$ pour la topologie dont un système fondamental de voisinages
de 0 est formé des $S^{-1}\mathfrak J_\lambda$.
\end{proposition}

En effet, si $v_\lambda$ est l'homomorphisme canonique
$S^{-1}A\to S_\lambda^{-1}A_\lambda$ déduit de $u_\lambda$, le noyau de
$v_\lambda$ est $S^{-1}\mathfrak J_\lambda$ et $v_\lambda$ est surjectif,
d'où la proposition (7.2.1).

\begin{corollary}[7.6.3]
\label{0.7.6.3-fr}
Si $S'$ est l'image canonique de $S$ dans le séparé complété $\widehat A$ de
$A$, $A\{S^{-1}\}$ s'identifie canoniquement à
$\widehat A\{{S'}^{-1}\}$.
\end{corollary}

On notera que même si $A$ est séparé et complet, il n'en est pas de même de
$S^{-1}A$ pour la topologie définie par les
$S^{-1}\mathfrak J_\lambda$, comme on le voit par exemple en prenant pour $S$
l'ensemble des $f^n$ ($n\geq0$), où $f$ est topologiquement nilpotent et non
nilpotent~: en effet, $S^{-1}A$ n'est pas réduit à 0 et d'autre part, pour tout
$\lambda$ il existe $n$ tel que $f^n\in\mathfrak J_\lambda$, donc
$1=f^n/f^n\in S^{-1}\mathfrak J_\lambda$ et
$S^{-1}\mathfrak J_\lambda=S^{-1}A$.

